% Part: proof-theory
% Chapter: proof-search
% Section: search-algorithm

\documentclass[../../../include/open-logic-section]{subfiles}

\begin{document}

\olfileid{pt}{ps}{sal}

\olsection{الگوریتم جست‌وجو برای \Log{G3c}}

الگوریتمی برای جست‌وجوی برهان در \Log{G3c} توصیف می‌کنیم. این تنها
الگوریتم ممکن و حتی کارآمدترین الگوریتم نیست، اما درست است: اگر
دنبالهٔ $\Gamma \Sequent \Delta$ دارای !!a{proof} باشد، الگوریتم
سرانجام متوقف می‌شود و !!a{proof}ی می‌یابد. همچنین خطاناپذیر است: اگر
بی‌آنکه !!a{proof}ی یافته باشد متوقف شود، یا هرگز متوقف نشود، آنگاه
$\Gamma \Sequent \Delta$ از آغاز معتبر نبوده است.

ویژگی ضروری الگوریتم این است که تضمین کند هر !!{formula} در
$\Gamma \Sequent \Delta$، یا در دنباله‌ای که هنگام جست‌وجوی برهان
تولید می‌شود، «کاهش» یابد؛ یعنی به‌عنوان فرمول اصلی قاعده‌ای استنتاجی
که وارونه اعمال می‌شود به کار رود، و این کار هرچند بار که لازم است
انجام گیرد. این ویژگی را \emph{انصاف} می‌نامیم.

الگوریتم مرحله‌به‌مرحله عمل می‌کند. خروجی هر مرحله درختی از دنباله‌هاست
و خروجی مرحلهٔ بعد با کاهش !!a{formula} در هر دنبالهٔ بالایی که از پیش
اصل نیست تولید می‌شود.

در مرحلهٔ~$0$ با $\Gamma \Sequent \Delta$ آغاز می‌کنیم. برای تضمین
انصاف، عددهایی را همچون اندیس به !!{formula}s موجود در $\Gamma$ و $\Delta$
اختصاص می‌دهیم، به‌گونه‌ای که !!{formula}s متفاوت عددهای متفاوت
بگیرند. برای نمونه، می‌توان !!{formula}s را از چپ به راست شمرد. عددها را
به‌صورت بالانویس می‌نویسیم. مثلاً اگر در پی !!a{proof}ی برای
$!A, !B \Sequent !C$ باشیم، خروجی مرحلهٔ~$0$ برابر
$!A^1, !B^2 \Sequent !C^3$ است. بزرگ‌ترین~$n$ که در چنین دنبالهٔ
اندیس‌داری به‌صورت بالانویس رخ دهد، \emph{اندیس بیشینهٔ} آن است
(در مثال،~$3$).

اگر دنباله سور داشته باشد، باید بدانیم هنگام کاهش
$\lexists[x][!A(x)]$ در راست یا $\lforall[x][!A(x)]$ در چپ از چه
ترم‌هایی استفاده کنیم. تنها به ترم‌هایی نیاز داریم که از
!!{constant}s $\Obj c_0$، $\Obj c_1$، $\Obj c_2$، \dots، و نمادهای
تابعی رخ‌داده در~$\Gamma$ و $\Delta$ ساخته می‌شوند. فرض می‌کنیم مجموعهٔ
همهٔ این ترم‌ها به ترتیبی ثابت شمارش شده باشد، تا بتوانیم مثلاً از
«نخستین !!{constant} که در $\Gamma \Sequent \Delta$ رخ نمی‌دهد» سخن
بگوییم. هر دنبالهٔ اندیس‌دار در درخت تولیدی الگوریتم، مجموعه‌ای متناظر
از !!{constant}s دارد. مجموعهٔ !!{constant}s اختصاص‌یافته به دنباله
در مرحلهٔ~$0$ مجموعهٔ همهٔ !!{constant}s رخ‌داده در آن است، اگر
چنین ثابتی وجود داشته باشد؛ و اگر هیچ‌یک از !!{constant}s را نداشته باشد،
$\{\Obj c_0\}$ است.

فرض کنید الگوریتم در مرحلهٔ~$n$ درختی از دنباله‌های اندیس‌دار با
مجموعه‌های متناظر !!{constant}s تولید کرده است. در مرحلهٔ~$n+1$ برای هر
دنبالهٔ بالایی کار زیر را انجام می‌دهیم.

اگر !!a{formula}~$!A$ هم در چپ و هم در راست دنباله رخ دهد، یا سمت چپ
شامل~$\lfalse$ باشد، کاری نمی‌کنیم: چنین دنباله‌ای در~\Log{G3c}
دارای !!{proof}s است و برای این دنبالهٔ بالایی خاص کاری باقی نمانده
است.

اگر هیچ !!{formula} غیراتمی‌ای در دنباله نباشد، الگوریتم را متوقف
می‌کنیم. برای $\Gamma \Sequent \Delta$ برهانی وجود ندارد.

در غیر این صورت، !!{formula} غیراتمی~$!A$ با کوچک‌ترین اندیس~$i$ را
برمی‌گزینیم. دنبالهٔ بالایی موردنظر یا
$!A^i, \Pi \Sequent \Lambda$ است یا
$\Pi \Sequent \Lambda, !A^i$. بگذارید $k$ اندیس بیشینهٔ دنباله باشد
و $C$ مجموعهٔ !!{constant}s اختصاص‌یافته به آن باشد.

~‌$!A^i$ را «کاهش» می‌دهیم؛ یعنی دنباله‌های بالایی تازه‌ای را برحسب
قاعدهٔ استنتاجی تعیین‌شده به‌وسیلهٔ !!{main operator}~$!A$ و برحسب
اینکه $!A^i$ در چپ یا راست ظاهر شده است، تولید می‌کنیم.

\begin{enumerate}
  \item اگر $!A \ident !B \land !C$ و در چپ رخ دهد، یک دنبالهٔ
  بالایی تازه بالای $(!B \land !C)^i, \Pi \Sequent \Lambda$ می‌افزاییم:
  \[
  \Axiom$!B^{k+1}, !C^{k+2}, \Pi \fCenter \Lambda$
  \RightLabel{\LeftR{\land}}
  \UnaryInf$(!B \land !C)^i, \Pi \fCenter \Lambda$
  \DisplayProof
  \]
  مجموعهٔ !!{constant}s اختصاص‌یافته به دنبالهٔ تازه~$C$ است.
  \item اگر $!A \ident !B \land !C$ و در راست رخ دهد، دو دنبالهٔ
  بالایی تازه بالای $\Pi \Sequent \Lambda, (!B \land !C)^i$ می‌افزاییم:
  \[
  \Axiom$\Pi \fCenter \Lambda, !B^{k+1}$
  \Axiom$\Pi \fCenter \Lambda, !C^{k+1}$
  \RightLabel{\RightR{\land}}
  \BinaryInf$\Pi \fCenter \Lambda, (!B \land !C)^i$
  \DisplayProof
  \]
  مجموعهٔ !!{constant}s اختصاص‌یافته به دنباله‌های تازه~$C$ است.

  \item حالت‌های $!A \ident \lnot !B$، $!A \ident !B \lor !C$ و
  $!A \ident !B \lif !C$ مشابه‌اند و به‌عنوان تمرین واگذار می‌شوند.

  \item اگر $!A \ident \lexists[x][!B(x)]$ و در چپ رخ دهد، یک دنبالهٔ
  بالایی تازه بالای $\lexists[x][!B(x)]^i, \Pi \Sequent \Lambda$
  می‌افزاییم:
  \[
  \Axiom$!B(c)^{k+1}, \Pi \fCenter \Lambda$
  \RightLabel{\LeftR{\lexists}}
  \UnaryInf$\lexists[x][!B(x)]^i, \Pi \fCenter \Lambda$
  \DisplayProof
  \]
  که در آن $c$ نخستین !!{constant} ناموجود در~$C$ است. مجموعهٔ
  !!{constant}s اختصاص‌یافته به دنبالهٔ تازه $C \cup \{c\}$ است.
  \item اگر $!A \ident \lexists[x][!B(x)]$ و در راست رخ دهد، یک
  دنبالهٔ بالایی تازه بالای
  $\Pi \Sequent \Lambda, \lexists[x][!B(x)]$ می‌افزاییم:
  \[
  \Axiom$\Pi \fCenter \Lambda, \lexists[x][!B(x)]^{k+1}, !B(t)^{k+2}$
  \RightLabel{\RightR{\lexists}}
  \UnaryInf$\Pi \fCenter \Lambda, \lexists[x][!B(x)]^i$
  \DisplayProof
  \]
  که در آن $t$ نخستین ترمی است که تنها شامل !!{constant}s موجود در~$C$ است و
  پیش‌تر در کاهش $\lexists[x][!B(x)]$ در سمت راست، زیر این دنباله، به
  کار نرفته است؛ و اگر همهٔ چنین ترم‌هایی به کار رفته باشند، نخستین
  ترم ساخته‌شده از~$C$ را دوباره به کار می‌بریم. مجموعهٔ
  !!{constant}s اختصاص‌یافته به دنبالهٔ تازه~$C$ است.
  \item حالت‌های $!A \ident \lforall[x][!B(x)]$ مشابه‌اند و به‌عنوان
  تمرین واگذار می‌شوند.
\end{enumerate}

اگر در مرحلهٔ $n+1$ بالای هیچ دنبالهٔ بالایی، دنبالهٔ تازه‌ای نیفزوده
باشیم، الگوریتم با موفقیت پایان می‌یابد. در این حالت !!a{proof}ی برای
~$\Gamma \Sequent \Delta$ یافته‌ایم که با پاک‌کردن همهٔ اندیس‌ها از
درخت مرحلهٔ $n+1$ حاصل می‌شود. همهٔ دنباله‌های بالایی به شکل
$!A, \Pi \Sequent \Lambda, !A$ یا
$\lfalse, \Pi \Sequent \Lambda$ هستند. همهٔ استنتاج‌ها استنتاج‌های
درست~\Log{G3c}اند. این امر برای استنتاج‌های گزاره‌ای آشکار است. توجه
کنید که در هر گام مجموعهٔ~$C$ از !!{constant}s اختصاص‌یافته به یک
دنباله، همهٔ !!{constant}s رخ‌داده در آن دنباله را در بر دارد.
بنابراین، در کاهش با \LeftR{\lexists} و \RightR{\lforall}، شرط متغیر
ویژه برقرار است: چون متغیر ویژهٔ~$c$ !!a{constant}ی نبود که در
مجموعهٔ~$C$ اختصاص‌یافته به نتیجه باشد، در نتیجه نیز رخ نمی‌داد. پس
داریم:

\begin{prop}
الگوریتم جست‌وجوی برهان برای \Log{G3c} درست است: اگر با موفقیت پایان
یابد، دنبالهٔ آغازین~$\Gamma \Sequent \Delta$ دارای !!a{proof} است.
\end{prop}





\end{document}
